\section{ONDES}

% ============================================================
\subsection{Généralités}
\begin{formules}
  \ligne{Vitesse de propagation, célérité}
    {$v,\ c$}
    {m/s}
    {$v$ : vitesse de propagation d'une perturbation dans un milieu ; $c$ : vitesse de la lumière, vide : \newline
     \bleu{$c = 299\,792\,458$ m/s}}
  \ligne{Équation d'onde (d'Alembert)}
    {$\begin{gathered}
      \text{Axe Ox : } \frac{\partial^2\xi}{\partial x^2} = \frac{1}{v^2}\,\frac{\partial^2\xi}{\partial t^2} \\
      \text{Espace : } \nabla^2\vec\xi = \frac{1}{v^2}\,\frac{\partial^2\vec\xi}{\partial t^2}
     \end{gathered}$}
    {$\frac{[\xi]}{\mathrm{m}^2}$}
    {$\xi(x,t)$ : grandeur physique portant l'onde ; $\vec\xi(x,t)$ : idem sous forme de champ vectoriel}
  \ligne{Solution générale}
    {$\begin{gathered}
      \text{Progressive (sens $+x$) :} \\
      \xi(x,t) = f(x - vt) \\
      \text{Rétrograde (sens $-x$) :} \\
      \xi(x,t) = f(x + vt)
     \end{gathered}$}
    {$[\xi]$}
    {$x$ : position [m] ; $t$ : temps [s] ; $\xi$ : grandeur physique portant l'onde}
  \ligne{Contrainte}
    {$\sigma(x,t) = E\,\dfrac{\partial\xi(x,t)}{\partial x}$}
    {Pa}
    {$E$ : module d'élasticité}
  \ligne{Additivité des ondes}
    {$\xi_r(x,t) = \sum\limits_{i=1}^{N}\xi_i(x,t)$}
    {$[\xi]$}
    {$N$ : nombre d'ondes de même nature en $x$ ; $\xi_i$ : onde $i = 1..N$ ; $\xi_r$ : onde résultante}
  \soustitre{Solution périodique sinusoïdale (oscillateur)}
  \ligne{Fonction d'onde (onde de déplacement)}
    {$\begin{gathered}
      \xi(x,t) = \xi_0\sin(kx - \omega t + \varphi) \\
      \cos(x) = \sin\left(\tfrac{\pi}{2} - x\right) \\
      \sin\left(\alpha + \tfrac{\pi}{2}\right) = \cos(\alpha)
     \end{gathered}$}
    {$[\xi]$}
    {$\xi_0$ : amplitude $[\xi]$ ; $\varphi$ : phase [rad] \newline
     $v_\xi(x,t) = \dfrac{\partial}{\partial t}\,\xi(x,t)$}
  \ligne{Nombre d'onde}
    {$k = \dfrac{2\pi}{\lambda} = \dfrac{\omega}{v}$}
    {m$^{-1}$}
    {$T$ : période [s]}
  \ligne{Longueur d'onde}
    {$\begin{gathered}
      \lambda = vT = \frac{v}{f} \\
      380\,\text{nm} \le \lambda \le 780\,\text{nm}
     \end{gathered}$}
    {m}
    {$f$ : fréquence [Hz]}
  \ligne{Pulsation}
    {$\omega = \dfrac{2\pi}{T} = 2\pi f$}
    {rad/s}
    {$f = \dfrac{1}{T} \Leftrightarrow T = \dfrac{1}{f}$ \newline $v = \lambda f = \dfrac{\lambda}{T}$}
  \ligneschema{
    \begin{scope}
      \draw[->] (0,0) -- (2.4,0) node[right] {$x$};
      \draw[->] (0,-0.4) -- (0,0.5) node[left] {$\xi$};
      \draw[thick, domain=0:2.2, samples=60] plot (\x, {0.3*sin(deg(2*pi*\x/1.0))});
      \draw[<->] (0.25,0.42) -- (1.25,0.42) node[midway, above] {$\lambda$};
      \draw[->] (0.3,-0.48) -- (0.75,-0.48) node[midway, below] {$\vec v$};
      \node at (2.0,-0.32) {$t$ fixe};
    \end{scope}
    \begin{scope}[xshift=3.3cm]
      \draw[->] (0,0) -- (2.4,0) node[right] {$t$};
      \draw[->] (0,-0.4) -- (0,0.5) node[left] {$\xi$};
      \draw[thick, domain=0:2.2, samples=60] plot (\x, {-0.3*sin(deg(2*pi*\x/1.0))});
      \draw[<->] (0.75,0.42) -- (1.75,0.42) node[midway, above] {$T$};
      \node at (2.0,-0.32) {$x$ fixe};
    \end{scope}
  }
  \soustitre{Fréquences propres}
  \ligne{Fermé-fermé}
    {$f_n = \dfrac{nv}{2L}$}
    {Hz}
    {$v$ : vitesse de propagation ; $L$ : longueur de la cavité}
  \ligne{Fermé-ouvert}
    {$f_n = \dfrac{(2n-1)v}{4L}$}
    {Hz}
    {$n$ : entier positif (1, 2, 3...)}
  \ligne{Ouvert-ouvert}
    {$f_n = \dfrac{nv}{2L}$}
    {Hz}
    {$n = 1$ : mode fondamental ; \newline $n > 1$ : harmoniques}
  \ligne{Solution onde sinusoïdale}
    {$\begin{gathered}
      \text{Mode $n$ : } \xi_n(x,t) = \\
      A_n\sin\left(\tfrac{n\pi}{L}x\right)\sin(2\pi f_nt + \varphi_n) \\
      \text{Générale : } \xi_r(x,t) = \sum\limits_{n=1}^{\infty}\xi_n(x,t)
     \end{gathered}$}
    {$[\xi]$}
    {}
\end{formules}

% ============================================================
\subsection{Ondes stationnaires : fréquences propres}
\begin{formules}
  \ligne{Solution générale (séparation des variables)}
    {$\xi(x,t) = A\sin(kx + \phi)\sin(\omega t + \varphi)$}
    {$[\xi]$}
    {Cavité 1D sur Ox. $L$ : longueur ; $A$ : amplitude ; $k$ : nombre d'onde ; $\omega$ : pulsation ;
     $\phi$ : phase spatiale ; $\varphi$ : phase temporelle}
  \ligne{Terme spatial (enveloppe)}
    {$\xi(x) = A\sin(kx + \phi)$}
    {$[\xi]$}
    {}
  \ligne{Conditions aux limites}
    {$\begin{gathered}
      \text{Fermé : } \xi = 0 \\
      \text{Ouvert : } \xi = \pm A
     \end{gathered}$}
    {}
    {Fermé : corde pincée, tube fermé... Ouvert : corde libre, bout de barreau, tube ouvert...}
  \ligne{Fermé-fermé}
    {$\begin{gathered}
      \xi(0) = 0 \Rightarrow \sin\phi = 0 \Rightarrow \phi = 0 \\
      \xi(L) = 0 \Rightarrow \sin(kL) = 0 \\
      \Rightarrow kL = n\pi \Rightarrow k_n = \tfrac{n\pi}{L}, \\
      f = \tfrac{vk}{2\pi} = \tfrac{v}{2\pi}\tfrac{n\pi}{L} \\
      \boxed{f_n = \frac{nv}{2L}} \\
      \boxed{\xi_n(x) = A\sin\left(\frac{n\pi}{L}x\right)}
     \end{gathered}$}
    {Hz}
    {Corde de guitare, barreau solide, boîte fermée... \newline
     $n = 1,2,3\dots$ ; $0 \le x \le L$ \newline
     $v = \lambda f$ et $k = \frac{2\pi}{\lambda}$}
  \ligne{Fermé-ouvert}
    {$\begin{gathered}
      \xi(0) = 0 \Rightarrow \phi = 0 \\
      \xi(L) = \pm A \Rightarrow \sin(kL) = \pm1 \\
      \Rightarrow kL = n\pi - \tfrac{\pi}{2} \\
      \Rightarrow k_n = \left(n - \tfrac{1}{2}\right)\tfrac{\pi}{L} \\
      f = \tfrac{vk}{2\pi} = \left(n - \tfrac{1}{2}\right)\tfrac{v}{2L} \\
      \boxed{f_n = \frac{(2n-1)v}{4L}} \\
      \boxed{\xi_n(x) = A\sin\left\{\left(n - \frac{1}{2}\right)\frac{\pi}{L}x\right\}}
     \end{gathered}$}
    {Hz}
    {Cavité avec ouverture, verre, bouteille, flûte de pan... \newline
     $n = 1,2,3\dots$ ; $0 \le x \le L$}
  \ligne{Ouvert-ouvert}
    {$\begin{gathered}
      \xi(0) = \pm A \Rightarrow \sin\phi = \pm1 \Rightarrow \\
      \phi = \pm\tfrac{\pi}{2} \\
      \xi(L) = \pm A \Rightarrow \cos(kL) = \pm1 \\
      \Rightarrow kL = n\pi \Rightarrow k_n = \tfrac{n\pi}{L} \\
      \boxed{f_n = \frac{nv}{2L}} \\
      \boxed{\xi_n(x) = A\cos\left(\frac{n\pi}{L}x\right)}
     \end{gathered}$}
    {Hz}
    {Tube ouvert, flûte à bec, barreau solide... \newline
     Même $f_n$ qu'en fermé-fermé ; \newline $0 \le x \le L$}
\end{formules}

% ============================================================
\subsection{Ondes sinusoïdales : vitesse, impédance, puissance, intensité}
{\par\scriptsize\noindent
\renewcommand{\arraystretch}{1.5}%
\setlength{\tabcolsep}{1pt}%
\setlength{\lI}{0.13\columnwidth}%
\setlength{\lT}{\dimexpr(\columnwidth-\lI-14\tabcolsep-8\arrayrulewidth)/6\relax}%
\begin{tabular}{|>{\raggedright\arraybackslash}m{\lI}|*{6}{>{\centering\arraybackslash}m{\lT}|}}
  \hline
  & Milieu solide & Milieu liquide & Gaz parfait & Ressort & Corde & Onde électro\-magnétique dans le vide \\ \hline
  & OL & OL & OL & OL & OT & OT \\ \hline
  Rappel $\alpha$ & $E$ & $K$ & $\gamma p_a$ & $k_rL$ & $F$ & $1/\varepsilon_0$ \\ \hline
  Inertie $\beta$ & $\rho$ & $\rho$ & $M$ & $\mu$ & $\mu$ & $\mu_0$ \\ \hline
  Impédance \newline $\boldsymbol{\mathcal{Z} = \sqrt{\alpha\beta}}$ & $\sqrt{E\rho}$ & $\sqrt{K\rho}$ & $p_a\sqrt{\frac{\gamma M}{RT}}$ &
    $\sqrt{k_rL\mu}$ & $\sqrt{F\mu}$ & $\sqrt{\mu_0/\varepsilon_0}$ \newline \bleu{$= 376.8\,\Omega$} \\ \hline
  Propagation \newline $\boldsymbol{v = \sqrt{\alpha/\beta}}$ & $\sqrt{\frac{E}{\rho}} = \frac{\mathcal{Z}}{\rho}$ & $\sqrt{\frac{K}{\rho}} = \frac{\mathcal{Z}}{\rho}$ &
    $\sqrt{\frac{\gamma RT}{M}} = \frac{\mathcal{Z}}{\rho}$ & $\sqrt{\frac{k_rL}{\mu}} = \frac{\mathcal{Z}}{\mu}$ &
    $\sqrt{\frac{F}{\mu}} = \frac{\mathcal{Z}}{\mu}$ & $c = \frac{\mathcal{Z}}{\mu_0}$ \newline $c = \frac{1}{\sqrt{\varepsilon_0\mu_0}}$ \\ \hline
  Amplitude & $\xi_0$ ou \newline $\sigma_0 = \mathcal{Z}\omega\xi_0$ & $\xi_0$ & $\xi_0$ ou \newline $p_0 = \mathcal{Z}\omega\xi_0$ &
    $\xi_0$ & $\xi_0$ & $E_0$ ou $B_0$ \newline \bleu{$E_0 = c\cdot B_0$} \\ \hline
  Puissance \newline $\bar P$ & $\frac{1}{2}S\mathcal{Z}(\omega\xi_0)^2$ & $\frac{1}{2}S\mathcal{Z}(\omega\xi_0)^2$ & $\frac{1}{2\mathcal{Z}}Sp_0^{\,2}$ &
    $\frac{1}{2}\mathcal{Z}(\omega\xi_0)^2$ & $\frac{1}{2}\mathcal{Z}(\omega\xi_0)^2$ & $\frac{1}{2\mathcal{Z}}SE_0^{\,2}$ \\ \hline
  Intensité \newline $\bar I$ & $\frac{1}{2}\mathcal{Z}(\omega\xi_0)^2$ & $\frac{1}{2}\mathcal{Z}(\omega\xi_0)^2$ & $\frac{1}{2\mathcal{Z}}p_0^{\,2}$ &
    - & - & $\frac{1}{2\mathcal{Z}}E_0^{\,2}$ \\ \hline
  Energie \newline $\bar w$ & $\frac{1}{2}\rho(\omega\xi_0)^2$ & $\frac{1}{2}\rho(\omega\xi_0)^2$ & $\frac{1}{2}\rho\left(\frac{p_0}{\mathcal{Z}}\right)^2$ &
    \multicolumn{2}{>{\centering\arraybackslash}m{\dimexpr2\lT+2\tabcolsep+\arrayrulewidth\relax}|}{$\frac{1}{2}\mu(\omega\xi_0)^2\ \left[\frac{\mathrm{J}}{\mathrm{m}}\right]$} & $\varepsilon_0E_0^{\,2}$ \\ \hline
\end{tabular}\par}
{\footnotesize OL, OT : onde longitudinale, transversale. \hfill
$\langle W_E\rangle = \overline{W_E}$ = densité d'énergie moyenne\par}

\smallskip
\noindent\begin{minipage}{\columnwidth}
{\footnotesize\textbf{Paramètres}\par}
{\scriptsize\noindent
\begin{minipage}[t]{0.49\columnwidth}\raggedright
  $\alpha$ : facteur de rétablissement \\
  $\beta$ : facteur d'inertie \\
  $v$ : vitesse de propagation de l'onde [m/s] \\
  $\mathcal{Z}$ : impédance de l'onde [$\Omega$] \\
  $\bar P$ : puissance moyenne [W] \\
  $\bar I$ : intensité moyenne [W/m$^2$] \quad ($\bar I = \bar P/S$) \\
  $\bar w$ : densité d'énergie moyenne [J/m$^3$] \\
  \quad systèmes linéiques : $[\bar w]$ = J/m \\[4pt]
  $E$ : module d'élasticité du solide [Pa] \\
  $K$ : module de compression du liquide \ [Pa] \\
  $\rho$ : masse volumique [kg/m$^3$] \\
  $\mu$ : masse linéique [kg/m] \\
  $k_r$ : constante de rappel du ressort [N/m] \\
  $L$ : longueur du ressort [m] \\
  $F$ : tension exercée sur la corde [N] \\
  $T$ : température du gaz [K]
\end{minipage}\hfill
\begin{minipage}[t]{0.49\columnwidth}\raggedright
  $\gamma$ : exposant adiabatique \\
  \quad \bleu{$\gamma = 5/3$ monoatomique} \\
  \quad \bleu{$\gamma = 7/5$ diatomique} \\
  \quad \bleu{$\gamma = 4/3$ triatomique ou plus} \\
  $p_a$ : pression du gaz [Pa] \\
  $M$ : masse molaire du gaz [kg/mol] \\
  \qquad Air sec : \bleu{$\boldsymbol{M = 0.028\,964\,4}$ $\frac{\mathrm{kg}}{\mathrm{mol}}$} \\
  $R$ : constante des gaz parfaits [J/(mol$\cdot$K)] \\
  \qquad \bleu{$\boldsymbol{R = 8.314\,462\,618\,153\,24}$ $\frac{\mathrm{J}}{\mathrm{mol\cdot K}}$} \\
  $\xi_0$ : déplacement des atomes, des molécules [m] \\
  $\sigma_0$ : contrainte [Pa] \\
  $p_0$ : pression acoustique [Pa] \\
  $S$ : surface du front d'onde [m$^2$] \\
  $\varepsilon_0$ : permittivité du vide [As/(Vm)] \\
  $\mu_0$ : perméabilité du vide [VS/(Am)] $= 4\pi\cdot10^{-7}$ \\
  $c$ : vitesse de la lumière dans le vide [m/s]
\end{minipage}\par}
\end{minipage}\par\medskip

\begin{formules}[sans]
  \ligne{Exposant adiabatique}
    {$\begin{gathered}
      \gamma = \frac{C_p}{C_v} = \frac{\chi_T}{\chi_S} \\
      \gamma = \frac{i+2}{i}
     \end{gathered}$}
    {[--]}
    {$i = 3$ monoatomique : \bleu{$\gamma = 5/3$} \newline
     $i = 5$ diatomique : \bleu{$\gamma = 7/5$} \newline
     $i = 6$ triatomique ou plus : \bleu{$\gamma = 4/3$}}
  \ligne{Module de compression}
    {$K = \dfrac{1}{\chi}$}
    {Pa}
    {$\chi$ : compressibilité}
  \ligne{Masse linéique}
    {$\mu = \dfrac{m}{L}$}
    {kg/m}
    {}
  \ligne{CNTP}
    {$\begin{gathered}
      T = 0\,^\circ\text{C} = 273\,\text{K} \\
      P = 1013\cdot10^2\,\text{Pa}
     \end{gathered}$}
    {}
    {}
\end{formules}

% ============================================================
\subsection{Réflexion -- transmission}
\begin{formules}
  \soustitre{Amplitude}
  \ligne{Facteur de réflexion}
    {$\begin{gathered}
      \text{Méca : } R = \frac{\xi_{0r}}{\xi_{0i}} = \frac{\mathcal{Z}_1 - \mathcal{Z}_2}{\mathcal{Z}_1 + \mathcal{Z}_2} \\
      \text{OEM : } R = \frac{E_{0r}}{E_{0i}} = \frac{\mathcal{Z}_2 - \mathcal{Z}_1}{\mathcal{Z}_1 + \mathcal{Z}_2}
     \end{gathered}$}
    {1}
    {Passage d'un milieu 1 à un milieu 2 \newline
     $\xi_{0i}, E_{0i}$ : amplitude incidente ; $\xi_{0r}, E_{0r}$ : réfléchie ; $\xi_{0t}, E_{0t}$ : transmise}
  \ligne{Facteur de transmission}
    {$\begin{gathered}
      \text{Méca : } T = \frac{\xi_{0t}}{\xi_{0i}} = \frac{2\mathcal{Z}_1}{\mathcal{Z}_1 + \mathcal{Z}_2} \\
      \text{OEM : } T = \frac{E_{0t}}{E_{0i}} = \frac{2\mathcal{Z}_2}{\mathcal{Z}_1 + \mathcal{Z}_2}
     \end{gathered}$}
    {1}
    {$\mathcal{Z}_1, \mathcal{Z}_2$ : impédances des milieux 1 et 2 \newline
     $1 + R = T$ ; $T > 0$ ; \newline $-1 < R < 1$}
  \soustitre{Intensité, puissance}
  \ligne{Facteur de réflexion}
    {$R_e = \left(\dfrac{\mathcal{Z}_1 - \mathcal{Z}_2}{\mathcal{Z}_1 + \mathcal{Z}_2}\right)^2$}
    {1}
    {$R_e + T_e = 1$}
  \ligne{Facteur de transmission}
    {$T_e = \dfrac{4\mathcal{Z}_1\mathcal{Z}_2}{(\mathcal{Z}_1 + \mathcal{Z}_2)^2}$}
    {1}
    {$R_e, T_e > 0$}
\end{formules}

% ============================================================
\subsection{Effet Doppler}
\begin{formules}
  \ligne{Fréquences apparentes}
    {$\begin{gathered}
      \text{Ondes mécaniques :} \\
      f_{obs} = f_{em}\,\frac{v - v_{obs}}{v - v_{em}} \\
      \text{Ondes électromagnétiques :} \\
      f_{obs} = f_{em}\sqrt{\frac{c - V}{c + V}} \\
      = f_{em}\sqrt{\frac{1 - \beta}{1 + \beta}} \\
      \beta = \frac{V}{c} \\
      \frac{\lambda_{obs}}{\lambda_{em}} = \frac{f_{em}}{f_{obs}} = \sqrt{\frac{1 + \beta}{1 - \beta}}
     \end{gathered}$}
    {Hz}
    {$f_{obs}$ : fréquence perçue par l'observateur O ; $f_{em}$ : fréquence émise par l'émetteur E ;
     $\vec v$ : vitesse de l'onde ; $\vec v_{obs}$ : vitesse de l'observateur ; $\vec v_{em}$ : vitesse de l'émetteur ;
     $\vec V = \vec v_{obs} - \vec v_{em}$ : vitesse relative de l'observateur par rapport à l'émetteur ;
     $\vec c$ : vitesse de la lumière
     \schema{
      \draw[->, gray] (-0.1,0) -- (1.9,0);
      \fill (0,0) circle (1pt) node[below] {E};
      \fill (1.5,0) circle (1pt) node[below] {O};
      \draw[->, champ] (0,0.12) -- (0.45,0.12) node[midway, above] {$\vec v_{em}$};
      \draw[->, champ] (1.5,0.12) -- (1.9,0.12) node[midway, above] {$\vec v_{obs}$};
      \draw[->, blue!70!black] (0.55,0.35) -- (1.15,0.35) node[midway, above] {$\vec v, \vec c$};
     }
     $v, v_{obs}, v_{em}, V, c$ : projections sur l'axe EO}
  \soustitre{Doppler avec angles (ondes mécaniques)}
  \ligne{Fréquence perçue}
    {$f' = f\left(\dfrac{1 - \frac{v_o}{v}\cos(\theta_o)}{1 - \frac{v_s}{v}\cos(\theta_s)}\right)$}
    {Hz}
    {$f$ : fréquence émise par la source S ; $f'$ : fréquence perçue par l'observateur O ;
     $v$ : vitesse de propagation de l'onde ; $v_s, v_o$ : vitesses de la source et de l'observateur ;
     $\theta_s, \theta_o$ : angles entre $\vec v_s$, $\vec v_o$ et l'axe S$\to$O}
  \ligneschema{
    \begin{scope}[font=\scriptsize]
    \draw[thick] (-0.2,0) -- (5.6,0);
    \draw[->] (0,0) -- (1.3,0.6) node[midway, above left] {$\vec v_s$};
    \draw[densely dotted] (1.3,0.6) -- (1.3,0);
    \draw[->] (0.7,0) arc (0:25:0.7);
    \node at (0.95,0.15) {$\theta_s$};
    \node at (0,0) {$\ast$};
    \node[below] at (0,-0.08) {S$_1$};
    \node[below] at (1.3,-0.08) {S$_2'$};
    \node at (0,-0.65) {$T$};
    \fill (3.6,0) circle (1pt);
    \draw[->] (3.6,0) -- (4.8,0.85) node[midway, above left] {$\vec v_o$};
    \draw[densely dotted] (4.8,0.85) -- (4.8,0);
    \draw[->] (4.15,0) arc (0:35:0.55);
    \node at (4.45,0.17) {$\theta_o$};
    \node[below] at (3.6,-0.08) {O$_1$};
    \node[below] at (4.8,-0.08) {O$_2'$};
    \node at (3.6,-0.65) {$T'$};
    \end{scope}
  }
  \lignepleine{\textcolor{rose}{Exemples : $v = 340$ m/s (son dans l'air) ; $f = 400$ Hz ;
     $v_o = 1.4$ m/s (piéton) ; $v_s = 100$ km/h $\approx 30$ m/s (ambulance)}}
  \soustitre{Source immobile : $v_s = 0$ \quad $f' = f\left(1 - \frac{v_o}{v}\cos(\theta_o)\right)$}
  \ligne{Observateur va vers la source}
    {$\theta_o = \pi$ \newline
     $f' = f\left(1 + \dfrac{v_o}{v}\right)$}
    {Hz}
    {\schema{
      \draw[->] (0.9,0) -- (0,0) node[midway, below] {$\vec v_o$};
      \fill (0.9,0) circle (0.8pt) node[right] {O};
     }
     \textcolor{rose}{$= 401.65$ Hz}}
  \ligne{Observateur s'éloigne de la source}
    {$\theta_o = 0$ \newline
     $f' = f\left(1 - \dfrac{v_o}{v}\right) = f\left(\dfrac{v - v_o}{v}\right)$}
    {Hz}
    {\schema{
      \fill (0,0) circle (0.8pt) node[left] {O};
      \draw[->] (0,0) -- (0.9,0) node[midway, below] {$\vec v_o$};
     }
     \textcolor{rose}{$= 398.35$ Hz}}
  \soustitre{Observateur immobile : $v_o = 0$ \quad $f' = f\,\dfrac{1}{1 - \frac{v_s}{v}\cos(\theta_s)}$\rule[-1.6ex]{0pt}{4ex}}
  \ligne{Source se rapproche de l'observateur}
    {$\theta_s = 0$ \newline
     $f' = \dfrac{f}{1 - \frac{v_s}{v}} = f\left(\dfrac{v}{v - v_s}\right)$}
    {Hz}
    {\schema{
      \node at (0,0) {$\ast$}; \node[below] at (0,-0.05) {S};
      \draw[->] (0.05,0) -- (0.8,0) node[midway, above] {$\vec v_s$};
      \node at (1.4,0) {$\times$}; \node[below] at (1.4,-0.05) {O};
     }
     \textcolor{rose}{$= 438.7$ Hz}}
  \ligne{Source s'éloigne de l'observateur}
    {$\theta_s = \pi$ \newline
     $f' = \dfrac{f}{1 + \frac{v_s}{v}} = f\left(\dfrac{v}{v + v_s}\right)$}
    {Hz}
    {\schema{
      \draw[->] (0.75,0) -- (0,0) node[midway, above] {$\vec v_s$};
      \node at (0.8,0) {$\ast$}; \node[below] at (0.8,-0.05) {S};
      \node at (1.4,0) {$\times$}; \node[below] at (1.4,-0.05) {O};
     }
     \textcolor{rose}{$\approx 367.6$ Hz}}
\end{formules}

% ============================================================
\subsection{Décibel}
\begin{formules}
  \ligne{Intensité moyenne en décibel}
    {$I_{dB} = 10\log_{10}\dfrac{\bar I}{I_{ref}}$}
    {dB}
    {$\bar I$ : intensité moyenne ; $I_{ref}$ : intensité de référence. dB acoustique : \newline
     $I_{ref} = 10^{-12}$ W/m$^2$}
\end{formules}

% ============================================================
\subsection{Générale ondes}
\begin{formules}
  \ligne{Intensité instantanée}
    {$I(x,t) = \dfrac{\partial P(x,t)}{\partial S}$}
    {W/m$^2$}
    {}
  \ligne{Puissance instantanée}
    {$P(x,t) = \displaystyle\int_S \vec I(x,t)\cdot d\vec S$}
    {W}
    {Si $I$ uniforme et $\perp$ à $S$ : \newline $P(x,t) = S\cdot I(x,t)$}
\end{formules}

% ============================================================
\subsection{Ondes mécaniques : barreau}
\begin{formules}
  \ligne{Puissance}
    {$\begin{gathered}
      P(x,t) = \\
      S\sqrt{E\rho}\,\omega^2\xi_0^2\cos^2(kx - \omega t)
     \end{gathered}$}
    {W}
    {$S$ : section du barreau}
  \ligne{Puissance moyenne (onde plane élastique)}
    {$\bar P = \dfrac{S}{2}\sqrt{E\rho}\,\omega^2\xi_0^2$}
    {W}
    {}
  \ligne{Intensité instantanée}
    {$\begin{gathered}
      I(x,t) = \\
      \sqrt{E\rho}\,\omega^2\xi_0^2\cos^2(kx - \omega t) \\
      I = \frac{P}{S}
     \end{gathered}$}
    {W/m$^2$}
    {}
  \ligne{Intensité moyenne}
    {$\bar I = \dfrac{1}{2}\sqrt{E\rho}\,\omega^2\xi_0^2 = \dfrac{\bar P}{S}$}
    {W/m$^2$}
    {}
  \ligne{Densité d'énergie instantanée}
    {$w = \xi_0^2\,\rho\,\omega^2\cos^2(kx - \omega t)$}
    {J/m$^3$}
    {}
  \ligne{Densité d'énergie moyenne}
    {$\bar w = \dfrac{1}{2}\rho\,\omega^2\xi_0^2$}
    {J/m$^3$}
    {}
  \ligne{Puissance onde plane}
    {$\begin{gathered}
      P(x,t) = \\
      ESk\omega\xi_0^2\cos^2(kx - \omega t)
     \end{gathered}$}
    {W}
    {}
\end{formules}

% ============================================================
\subsection{Ondes sonores}
\begin{formules}
  \ligne{Intensité de référence}
    {$I_{ref} = 10^{-12}$ W/m$^2$}
    {W/m$^2$}
    {Seuil d'audition. Seuil de douleur : $I_d = 1$ W/m$^2$}
  \ligne{Intensité en dB}
    {$I_{dB} = 10\log_{10}\dfrac{\bar I}{I_{ref}}$}
    {dB}
    {}
  \ligne{Surpression}
    {$\begin{gathered}
      p(x,t) = p_0\sin(kx - \omega t) \\
      = -\rho v^2\,\frac{\partial\xi(x,t)}{\partial x}
     \end{gathered}$}
    {Pa}
    {}
  \ligne{Pression dynamique (amplitude de pression)}
    {$\begin{gathered}
      p_0 = \sqrt{2\mathcal{Z}\bar I} = \rho v\omega\xi_0 \\
      = \mathcal{Z}\omega\xi_0 = \rho v^2k\xi_0
     \end{gathered}$}
    {Pa}
    {Pour une onde progressive, la pression acoustique est déphasée de $\pi/2$ par rapport au déplacement}
  \ligne{Onde sonore plane}
    {$\begin{gathered}
      P_{tot}(x,t) = \\
      P_{atm} + p_0\sin(kx - \omega t + \phi)
     \end{gathered}$}
    {Pa}
    {$P_{tot}$ : pression totale}
  \ligne{Onde sonore sphérique}
    {$\begin{gathered}
      P_{tot}(r,t) = \\
      P + \sqrt{\frac{\mathcal{Z}\bar P_0}{2\pi}}\,\frac{\sin(kr - \omega t + \phi)}{r}
     \end{gathered}$}
    {Pa}
    {$\bar P_0$ : puissance moyenne de la source}
  \ligne{Intensité moyenne (onde acoustique plane)}
    {$\bar I = \dfrac{1}{2}\rho v\omega^2\xi_0^2 = \dfrac{\bar P}{S}$}
    {W/m$^2$}
    {Source sphérique : $S = 4\pi r^2$}
  \ligne{Densité d'énergie moyenne}
    {$\bar w = \dfrac{1}{2}\rho\,\omega^2\xi_0^2 = \dfrac{\bar I}{v}$}
    {J/m$^3$}
    {}
\end{formules}

% ============================================================
\subsection{Onde électromagnétique}
\begin{formules}
  \ligne{Densités d'énergie moyennes}
    {$\begin{gathered}
      \bar w_E = \frac{1}{4}\varepsilon_0E^2 \quad \bar w_B = \frac{B^2}{4\mu_0} \\
      \bar w = \bar w_E + \bar w_B
     \end{gathered}$}
    {J/m$^3$}
    {$E, B$ : amplitudes des champs}
  \ligne{Densités d'énergie instantanées}
    {$\begin{gathered}
      w_E = \frac{1}{2}\varepsilon_0E^2 \quad w_B = \frac{B^2}{2\mu_0} \\
      w = w_E + w_B
     \end{gathered}$}
    {J/m$^3$}
    {}
  \ligne{Pression de radiation}
    {$\begin{gathered}
      \text{Isotrope : } P_{rad} = \tfrac{1}{3}\chi w \\
      \text{Plan : } P_{rad} = \chi w \\
      \text{Force : } F = S\cdot P_{rad}
     \end{gathered}$}
    {Pa \newline\newline N}
    {Isotrope : toutes directions ; plan : une direction ; $F$ : force transmise par l'onde à la surface.
     $\chi$ : coefficient lié à la surface : tout absorbé $= 1$, tout réfléchi $= 2$, absorption/réflexion $= 1.5$}
  \ligne{Vecteur de Poynting}
    {$\vec R = \dfrac{\vec E\wedge\vec B}{\mu_0} \quad \bar I = \dfrac{1}{2}R_0$}
    {W/m$^2$}
    {$\vec E\perp\vec B \to$ onde électromagnétique ; $R_0$ : amplitude de $R$}
\end{formules}
